\section{Evidence: a six-year Veneto panel}\label{sec:evidence}

\subsection{Data}

The panel is drawn from the Veneto Worker History data that KSS use, for the years 1996 to 2001:
\R{real.n.person_years} person-years of \R{real.n.workers} workers at \R{real.n.firms} firms, in the set of
firms that remains connected when any one worker is left out. Log wages are net of year effects and two age
terms, removed in a first step (Section~\ref{sec:controls}). \R{real.n.stayers} workers are stayers, observed for two
to six years (\R{real.n.stayer_rows} person-years, \R{real.share.stayer_py} of the total); the \R{real.n.movers}
movers have \R{real.n.mover_rows} matches. The design therefore has \R{real.n.rows} rows, \R{real.share.stayer_rows}
of them stayers' person-years. The numbers below come from a prototype that follows FEVC's algorithm, with
exact leverages for stayers and 1{,}600 projections for movers.

\subsection{Estimates and standard errors}

\begin{table}[H]
\centering\small
\caption{Six-year Veneto panel: estimates and standard errors}\label{tab:real}
\begin{tabular}{@{}lcccccc@{}}
\toprule
 & Plug-in & Leave-out & \multicolumn{2}{c}{Standard error} & First term, & Movers, \% \\
\cmidrule(lr){4-5}
Component & estimate & estimate $\hat\theta$ & $\times10^{-4}$ & \% of $\hat\theta$ & \% of $\hat\V[\hat\theta]$ & of first term \\
\midrule
\tabinput{gen/tab_real.tex}
\bottomrule
\multicolumn{7}{@{}p{0.94\textwidth}@{}}{\rule{0pt}{2.6ex}Plug-in: $\hat\beta'A\hat\beta$. Standard error: the square
root of \eqref{eq:Vhat} with 16 draws for the double sum. ``First term'': $4\sum_i\hat m_i^{\,2}\st_i^2$; it
exceeds 100\% because the double sum is subtracted. The last column is the share of movers' rows in that term.}
\end{tabular}
\end{table}

Four things stand out. The standard errors are small, between \R{real.relse.min} and \R{real.relse.max} of the
estimates: with five million person-years, sampling error is not the main uncertainty in these numbers.
Movers' rows carry most of the first term: nearly all of it for $\sigma^2_\alpha$, $\sigma^2_\psi$ and
$\sigma_{\alpha,\psi}$, and \R{real.share.movers.sum} for the variance of the sum, to which stayers contribute
\R{real.share.stayers.sum}. Stayers' rows thus supply nearly a quarter of the first term for the sum, with
variance estimates that need no model; the rest of that term, and nearly all of it for the other three
components, rests on the model for movers. Two computational pieces are small: the adjustment for stayers is at
most \R{real.adjust_max} of the variance estimate, and other sets of draws for the double sum, or
a run with 64 draws, change no standard error by more than \R{real.run_maxdiff}. Finally, the tails of the
wage distribution matter for this standard error: the \R{real.capped.movers} of movers' rows whose person effect
lies beyond the winsorising limit carry \R{real.capped.first_term.sum} of the first term for the variance of the
sum.

\subsection{The model for movers in the data}

\begin{figure}[H]
\centering
\def\profalphastay{(-1.826,2.678) (-0.846,1.459) (-0.523,0.960) (-0.298,0.757) (-0.109,0.660) (0.070,0.628) (0.262,0.641) (0.499,0.662) (0.874,0.784) (1.939,1.184)}
\def\profalphafit{(-1.826,2.249) (-0.846,1.268) (-0.523,1.048) (-0.298,0.928) (-0.109,0.849) (0.070,0.789) (0.262,0.741) (0.499,0.705) (0.874,0.695) (1.939,1.057)}
\def\profalphamove{(-1.713,2.424) (-0.838,1.389) (-0.522,0.948) (-0.299,0.810) (-0.110,0.760) (0.068,0.730) (0.258,0.736) (0.495,0.800) (0.867,0.881) (1.926,1.081)}
\def\profpsistay{(-1.868,1.352) (-0.796,1.076) (-0.414,1.021) (-0.176,1.028) (0.009,0.952) (0.187,0.937) (0.360,0.976) (0.564,0.962) (0.894,0.848) (1.838,0.913)}
\def\profpsifit{(-1.868,1.323) (-0.796,1.055) (-0.414,1.027) (-0.176,1.015) (0.009,0.981) (0.187,0.975) (0.360,0.949) (0.564,0.942) (0.894,0.904) (1.838,0.888)}
\def\profpsimove{(-1.868,1.451) (-0.798,1.041) (-0.419,0.988) (-0.178,0.942) (0.014,0.882) (0.183,0.884) (0.359,0.866) (0.571,0.931) (0.899,0.955) (1.680,1.124)}
\def\figalphalo{(-1.683,17.859) (-0.838,5.629) (-0.524,3.553) (-0.300,2.909) (-0.111,2.704) (0.067,2.634) (0.257,2.890) (0.493,3.626) (0.863,5.183) (1.932,9.254)}
\def\figalphafit{(-1.683,15.960) (-0.838,6.384) (-0.524,4.410) (-0.300,3.438) (-0.111,2.889) (0.067,2.565) (0.257,2.419) (0.493,2.479) (0.863,3.139) (1.932,10.138)}
\def\figpsilo{(-1.957,9.523) (-0.798,5.408) (-0.421,4.579) (-0.178,4.275) (0.014,4.235) (0.182,3.768) (0.356,3.728) (0.570,3.977) (0.904,3.852) (1.680,4.436)}
\def\figpsifit{(-1.957,9.363) (-0.798,5.346) (-0.421,4.734) (-0.178,4.301) (0.014,4.021) (0.182,3.933) (0.356,3.806) (0.570,3.706) (0.904,3.635) (1.680,4.392)}

\begin{tikzpicture}
\begin{groupplot}[group style={group size=2 by 1, horizontal sep=1.3cm}, width=0.5\textwidth, height=0.36\textwidth,
  ymin=0, ymax=20, xmin=-2.4, xmax=2.4, ymajorgrids, grid style={gray!20},
  xlabel style={font=\footnotesize}, ylabel style={font=\footnotesize}, ticklabel style={font=\footnotesize},
  title style={font=\footnotesize}, legend style={font=\footnotesize, draw=none, fill=none, legend cell align=left}]
\nextgroupplot[title={By estimated person effect}, xlabel={standardised person effect (decile means)},
  ylabel={variance of a mover's row ($\times100$)}, legend style={at={(0.98,0.97)}, anchor=north east}]
\addlegendimage{only marks, mark=*, mark size=1.8pt, cmove}
\addlegendentry{leave-out estimates, average}
\addlegendimage{no marks, thick, cstay}
\addlegendentry{fitted by the model, average}
\addplot[cmove, only marks, mark=*, mark size=1.8pt, forget plot] coordinates {\figalphalo};
\addplot[cstay, thick, no marks, forget plot] coordinates {\figalphafit};
\nextgroupplot[title={By estimated firm effect}, xlabel={standardised firm effect (decile means)}]
\addplot[cmove, only marks, mark=*, mark size=1.8pt] coordinates {\figpsilo};
\addplot[cstay, thick, no marks] coordinates {\figpsifit};
\end{groupplot}
\end{tikzpicture}
\caption{Movers' rows by decile of the estimated effects: average leave-out estimate $\hat\sigma_i^2$ and average
fitted variance $\st_i^2$. Deciles are those of the stayers' estimated effects.}
\label{fig:profile}
\end{figure}

Figure~\ref{fig:profile} shows what the model works with and what it delivers. The averages of the leave-out
estimates vary a great deal. They are \R{dec.data.lo.bottom} times their average in the lowest decile of the
estimated person effect, \R{dec.data.lo.mid} times in the middle deciles and \R{dec.data.lo.top} times in the
highest, and they fall by more than half across the deciles of the firm effect. Stayers show the same pattern
at a lower level: the average $q_g$ (weighted by $T_g-1$) is \R{real.prof.alpha.stay.lo} times the stayers' overall average in the
lowest decile of the person effect, about \R{real.prof.alpha.stay.mid} times in the middle deciles and
\R{real.prof.alpha.stay.hi} times in the highest. The fitted variances follow the leave-out estimates closely
across the firm effect. Across the person effect they reproduce the U shape, but not exactly: they are
\R{dec.data.over.min} to \R{dec.data.over.max} above the leave-out averages in deciles 2 to 4 and
\R{dec.data.under.min} to \R{dec.data.under.max} below them in deciles 7 to 9. This fit, a quadratic in effects capped at three standard deviations, falls short there.

The scale of step~1 has a clear role in the data as well. The leave-out estimates average \R{real.lo.len1} for
one-year matches and \R{real.lo.longer.min} to \R{real.lo.longer.max} for longer ones. The variance of a mover's
yearly wages within a match is only \R{real.within_match_mean} on average: most of the error of a mover's row is
common to the years of the match, which is why nothing is assumed within a match and why the scale depends on
its length.

Can the decile averages of Figure~\ref{fig:profile} be trusted as a benchmark? Averaged over a fixed set of many
rows, the leave-out estimates estimate the average variance of those rows. Here the rows are sorted by an
estimated effect that depends on the same errors, which can distort the averages. Table~\ref{tab:deciles} adds
three simulated settings, where the truth is known. There the average leave-out estimate of a decile differs
from the average true variance by at most \R{dec.sim.gap.max}, upwards in the extreme deciles and about 5\%
downwards in the middle. The distortion in the data need not be the same.

\begin{table}[H]
\centering\small
\caption{Movers' rows by decile of the estimated person effect: average leave-out estimate $\hat\sigma_i^2$,
average fitted variance $\st_i^2$ and, in simulated data, average true variance ($\times100$)}\label{tab:deciles}
\begin{tabular}{@{}lcccccccccc@{}}
\toprule
 & \multicolumn{10}{c}{Decile of the estimated person effect (lowest to highest)} \\
\cmidrule(l){2-11}
 & 1 & 2 & 3 & 4 & 5 & 6 & 7 & 8 & 9 & 10 \\
\midrule
\tabinput{gen/tab_deciles.tex}
\bottomrule
\multicolumn{11}{@{}p{0.92\textwidth}@{}}{\rule{0pt}{2.6ex}Deciles are those of the stayers' estimated person
effects, so movers' rows are spread unevenly over them (\R{dec.rows.min} to \R{dec.rows.max} of the rows per decile).
Simulated data: averages over \R{dec.reps} draws on the Veneto design (Section~\ref{sec:sim}); across the draws the
decile averages of the leave-out estimates vary by at most \R{dec.sd.max} (standard deviation). ``True variance''
is the average true variance of the rows that fall in the decile.}
\end{tabular}
\end{table}

In the simulated settings the fitted variances can also be compared with the truth. In the highest decile
they exceed the true variances in all three settings (by \R{dec.N1.fit.top}, \R{dec.N6.fit.top} and
\R{dec.N10.fit.top}), and in the setting without any profile in the lowest decile as well: the fit inherits
part of the distortion, which is the mechanism described in Section~\ref{sec:transfer}. What the discrepancies of the first panel mean for the
standard error, that is, how large $b(\Delta)/\V[\hat\theta]$ is on this panel, the table cannot say, and the
data cannot say it either, because movers' true variances are not observed. Scenario~10 of the simulation below
gives movers exactly the decile profile of the first panel; there the answer can be computed.

\subsection{Is any stayer or direction dominant?}

Sample versions of the two quantities in Proposition~\ref{prop:conc} are computed with the estimated
variances. For a stayer the first is
$\hat\phi_g=\sum_{i\in g}\st_i^2\,(4\hat m_i^{\,2}+2\hat r_i)/\hat\V[\hat\theta]$, where $\hat r_i$ is the simulated
value of $\sum_\ell C_{i\ell}^2\st_\ell^2$ from 32 draws as in \eqref{eq:sim}; the same ratio is formed for each
mover's row. The second is $\hat\lambda_{\max}^2/\hat\V[\hat\theta]$, with $\hat\lambda_{\max}^2$ from 20 power
iterations on the matrix with entries $\st_iC_{i\ell}\st_\ell$. These are concentration ratios, not exact shares:
$\hat m_i^{\,2}$ overstates $m_i^2$, so the $\hat\phi_g$ do not add up to one. All are small. The largest
$\hat\phi_g$ is \R{real.max_stayer_share}, the largest ratio of a mover's row \R{real.max_mover_share}, and
$\hat\lambda_{\max}^2/\hat\V[\hat\theta]$ is between \R{real.lambda.min} and \R{real.lambda.max} for the four
components.

\subsection{Simulation}\label{sec:sim}

\paragraph{Design.} The effects fitted on the panel serve as the truth, and new errors are drawn for every
person-year, \R{sim.reps} times in each of ten scenarios (\R{sim.failures} draws failed). Each draw runs the
estimator exactly as on the data: effects, leverages of movers from the same 1{,}600 projections, the two steps
of the model for movers with estimated effects in the profile, and 16 draws for the double sum. The variance
of $\hat\theta$ is computed from \eqref{eq:V} for each scenario. Appendix~\ref{app:sim} defines the scenarios and
reports how each was realised. In scenarios 1 to 7 movers and stayers share one variance profile. Scenario~1
has none. In scenarios 2 to 5 the profile is quadratic in the standardised person effect, as the model's is,
with one difference: the model caps the standardised effects at three standard deviations and the scenarios do
not. Scenario~6 uses a step profile taken from the data, the stayers' average sample variance by decile of the
person effect, and scenario~7 a profile that is log-linear in the wage level. Scenarios 3 to 5 also violate A3
(fat tails, skewness). In scenarios 8 to 10 movers have a profile of their own. In 8 and 9 their variances rise
three times as fast with the distance from the mean wage as the stayers', or not at all. In both, the movers'
profile lies in the model's quadratic basis before capping; exact recovery by the two regressions requires in
addition the condition on the scale in Corollary~\ref{cor:transfer}. Scenario~10 is calibrated to the data and is
not within the model: movers follow the decile profile of their own leave-out estimates in the first row of
Table~\ref{tab:deciles}.

\begin{table}[H]
\centering\small
\caption{Simulation on the Veneto design: square root of the mean variance estimate over the true variance}\label{tab:sim}
\begin{tabular}{@{}rlcccccccccc@{}}
\toprule
 & & \multicolumn{4}{c}{With the model for movers} & & \multicolumn{4}{c}{With movers' true variances} \\
\cmidrule(lr){3-6}\cmidrule(l){8-11}
 & Scenario & $\sigma^2_\alpha$ & $\sigma^2_\psi$ & $\sigma_{\alpha,\psi}$ & $\sigma^2_{\alpha+\psi}$ & & $\sigma^2_\alpha$ & $\sigma^2_\psi$ & $\sigma_{\alpha,\psi}$ & $\sigma^2_{\alpha+\psi}$ \\
\midrule
\tabinput{gen/tab_sim.tex}
\bottomrule
\multicolumn{11}{@{}p{0.9\textwidth}@{}}{\rule{0pt}{2.6ex}A value of 1 means that the mean variance estimate
equals the true variance; the entry is the root-mean-square standard error over the true one.
\R{sim.reps} draws per scenario. In scenarios 8 to 10 movers have a profile of their own. Right panel:
\eqref{eq:Vhat} with $\st_i^2=\sigma_i^2$ on movers' rows, which isolates the stayers' part.}
\end{tabular}
\end{table}

\begin{table}[H]
\centering\small\setlength{\tabcolsep}{4.6pt}
\caption{Simulation on the Veneto design: coverage of the 95\% interval}\label{tab:simcov}
\begin{tabular}{@{}rlcccccccccc@{}}
\toprule
 & & \multicolumn{4}{c}{With the model for movers} & & \multicolumn{4}{c}{With movers' true variances} \\
\cmidrule(lr){3-6}\cmidrule(l){8-11}
 & Scenario & $\sigma^2_\alpha$ & $\sigma^2_\psi$ & $\sigma_{\alpha,\psi}$ & $\sigma^2_{\alpha+\psi}$ & & $\sigma^2_\alpha$ & $\sigma^2_\psi$ & $\sigma_{\alpha,\psi}$ & $\sigma^2_{\alpha+\psi}$ \\
\midrule
\tabinput{gen/tab_simcov.tex}
\bottomrule
\multicolumn{11}{@{}p{0.9\textwidth}@{}}{\rule{0pt}{2.6ex}Share of draws in which
$\hat\theta\pm1.96\,\hat\V[\hat\theta]^{1/2}$ covers the true $\theta$. With \R{sim.reps} draws the Monte Carlo
standard error is about 0.015.}
\end{tabular}
\end{table}

\begin{table}[H]
\centering\small
\caption{The model's error at fixed characteristics, $[1+b(\Delta)/\V[\hat\theta]]^{1/2}$, against the simulated ratio}\label{tab:bdelta}
\begin{tabular}{@{}rlcccccccccc@{}}
\toprule
 & & \multicolumn{4}{c}{At fixed characteristics} & & \multicolumn{4}{c}{As implemented (Table~\ref{tab:sim})} \\
\cmidrule(lr){3-6}\cmidrule(l){8-11}
 & Scenario & $\sigma^2_\alpha$ & $\sigma^2_\psi$ & $\sigma_{\alpha,\psi}$ & $\sigma^2_{\alpha+\psi}$ & & $\sigma^2_\alpha$ & $\sigma^2_\psi$ & $\sigma_{\alpha,\psi}$ & $\sigma^2_{\alpha+\psi}$ \\
\midrule
\tabinput{gen/tab_bdelta.tex}
\bottomrule
\multicolumn{11}{@{}p{0.9\textwidth}@{}}{\rule{0pt}{2.6ex}Left: the two steps of the model are run with the true
variances as outcomes and the true effects as characteristics; this gives $v_i^*$, $\Delta_i=v_i^*-\sigma_i^2$ and
$b(\Delta)$ from \eqref{eq:b}, with no simulated outcomes. It concerns the model for movers only; the stayers' departures from A3 in scenarios 3 to 5 are not in it.
Right: the left panel of Table~\ref{tab:sim}. The panels differ in everything else: estimated effects, the
sampling error of the two fits, the stayers' part, projected weights, and the projection of the covariance
matrix described in the text.}
\end{tabular}
\end{table}

\paragraph{Findings.} \emph{The stayers' part.} With movers' true variances (right panels of
Tables~\ref{tab:sim} and~\ref{tab:simcov}) the ratio is between \R{sim.or.noskew.min} and \R{sim.or.noskew.max}
in the eight scenarios without skewness, including scenario~3 with fat tails. In the two scenarios with skewed
stayers' errors it is \R{sim.N4.or.sum} and \R{sim.N5.or.sum} for the variance of the sum: the shift that
Section~\ref{sec:results} predicts, with the predicted signs. The simulations thus agree with
Proposition~\ref{prop:mean} and the formulas for departures from A3, although those are for exact weights and
these runs use projected ones.

\emph{With the model for movers.} Across the \R{sim.n.cells} cells the ratio of Table~\ref{tab:sim}, the
root-mean-square standard error over the true one, is between \R{sim.all.min} and \R{sim.all.max}. It is below
0.95 in \R{sim.n.below95} cells and above 1.05 in \R{sim.n.above105} (counts here and below are before
rounding). For $\sigma^2_\alpha$ the ratio is between
\R{sim.comp.alpha.min} and \R{sim.comp.alpha.max}, for $\sigma^2_\psi$ between \R{sim.comp.psi.min} and
\R{sim.comp.psi.max}, for the covariance between \R{sim.comp.cov.min} and \R{sim.comp.cov.max}, and for the
variance of the sum between \R{sim.comp.sum.min} and \R{sim.comp.sum.max}. The ratio is thus below one
by \R{sim.under.max} at most (the variance of the sum in scenario~4, where the benchmark with movers' true
variances in the right panel is itself about three percent low), and above one by up to \R{sim.over.max}. The largest ratios are those for
$\sigma^2_\psi$ in the four scenarios with worker-specific variance factors, and for the sum in two of them. In
scenario~10, where movers have the decile profile of the data, the ratios are \R{sim.N10.alpha}, \R{sim.N10.psi},
\R{sim.N10.cov} and \R{sim.N10.sum}. The coverage of the 95\% interval is between \R{sim.all.cov.min} and
\R{sim.all.cov.max} (Table~\ref{tab:simcov}), against \R{sim.or.cov.min} to \R{sim.or.cov.max} with movers' true
variances. No draw failed. In no draw did the last floor bind for more than \R{sim.floor.final.max} of movers'
rows, and the floor on the scale never bound. The Monte Carlo standard error of the ratios is at most \R{sim.ratio.mcse.max}. One step of the prototype in
these runs is not in Section~\ref{sec:estimator}: before the variances were read off, negative eigenvalues of the
estimated covariance matrix of the three base components were set to zero. This can only raise a variance
estimate, and it applies to Tables~\ref{tab:sim} to~\ref{tab:bdelta} (right panel). On
the data that matrix is positive definite, with a smallest eigenvalue of \R{real8.eig.ratio} of the largest, so
the projection changes nothing there. For these simulated draws it was not recorded whether it changed
anything. The validation runs of Section~\ref{sec:validation} use the estimator without the projection, and on
the same design they give the same ratios to within \R{val.six.devdiff.max}.

\emph{The model at fixed characteristics.} Table~\ref{tab:bdelta} separates the model for movers from the
rest. Its left panel evaluates the quantity of A4 on this design: the error $b(\Delta)$ of the model when its
two steps see the true variances and the true effects. In scenarios 1 and 9 the model's error is zero to three decimals. In the other
eight, $[1+b(\Delta)/\V[\hat\theta]]^{1/2}$ is between \R{sim.bd.other.min} and \R{sim.bd.other.max} for the first
three components and between \R{sim.bd.sum.min} and \R{sim.bd.sum.max} for the variance of the sum: at fixed
characteristics the model would understate that standard error by \R{sim.bd.sum.short.min} to
\R{sim.bd.sum.short.max}. Two sources can be told apart. One is the cap on the standardised effects. In
scenarios 2 to 5 and 8 the true variances keep growing with the square of the effect beyond three standard
deviations, where the model holds its profile constant. Few rows lie there, \R{sim.cap.N4.rows3} of movers'
rows, but in scenario~4 they carry \R{sim.cap.N4.first3.whole} of the first term of \eqref{eq:V} for the variance
of the sum (\R{sim.cap.N4.first3} of the movers' part). In that scenario the entry of Table~\ref{tab:bdelta} for the sum is \R{sim.cap.N4.3.sum} with the cap at
3; with a cap at 4 it would be \R{sim.cap.N4.4.sum}, at 6 \R{sim.cap.N4.6.sum}, and without a cap
\R{sim.cap.N4.none.sum}. In scenario~8 the same four numbers are \R{sim.cap.N8.3.sum}, \R{sim.cap.N8.4.sum},
\R{sim.cap.N8.6.sum} and \R{sim.cap.N8.none.sum}. The cap protects the fit from a few rows with extreme estimated effects, and this is its
price when variances keep rising in the tails. The other source is easy to overlook. In scenarios 2, 3, 6 and
10 every worker has a variance factor of its own, which the matches of a mover share and which no
characteristic in the model captures. Such errors average out in the first part of $b(\Delta)$, in expectation
over the factors, but not in its second, the double sum, which is negative when the errors of the model are
positively related within a worker. For $\sigma^2_{\alpha+\psi}$ the double sum lowers the mean of the variance
estimate by \R{sim.bdquad.mix.min} to \R{sim.bdquad.mix.max} of the variance in these four scenarios, and by at
most \R{sim.bdquad.other.max} elsewhere (Table~\ref{tab:regime}); it contains all errors of the model, not only
those from the worker factors. A model that is right about the average variance of every group of rows can
therefore still give a standard error that is too small.

\emph{As implemented.} The procedure as implemented (right panel of Table~\ref{tab:bdelta}) lies above its
fixed-characteristics counterpart in \R{sim.bd.above.all} of the \R{sim.n.cells} cells; for the variance of the
sum the difference is between $\R{sim.bd.diff.sum.min}$ and $\R{sim.bd.diff.sum.max}$ (before rounding). The
ratios of Table~\ref{tab:sim} are therefore the net of two opposing contributions: a model error that lowers the
variance estimate, and an often larger net contribution of the implementation in the other direction. The two
panels differ in several ways at once (estimated effects in place of true ones, noisy outcomes in the two fits,
the stayers' part in scenarios 3 to 5, the projection of the covariance matrix), so the comparison does not
isolate a cause. In scenario~1 the model is exact at fixed characteristics, and as implemented the ratio for
the variance of the sum is \R{sim.N1.sum}: an upward discrepancy of the implementation with no model error
behind it. We have no result that predicts its sign in other designs; Section~\ref{sec:validation} looks at
two.

\emph{What the simulations do not show.} They use one design and ten scenarios of our choosing, and the form of
the model was chosen on the same scenarios (Appendix~\ref{app:sim}). They show how the standard error behaves when
movers' variances are as specified. They cannot show how large $b(\Delta)/\V[\hat\theta]$ is on the Veneto panel
itself, where movers' variances are not observed; scenario~10 is the closest we can come.

\subsection{Validation: fresh draws and two further designs}\label{sec:validation}

The simulations above served to choose the model. After that choice the estimator was frozen and run once
more, with acceptance thresholds written down before any task ran. There are three runs.
\begin{itemize}
\item \emph{Six-year panel, fresh draws.} The ten scenarios of Table~\ref{tab:sim} with new error draws.
\item \emph{A three-year window.} The years 1999 to 2001 of the same data: \R{val.y3.n.person_years} person-years
of \R{val.y3.n.workers} workers at \R{val.y3.n.firms} firms, \R{val.y3.n.rows} rows. Stayers have two or three
years (\R{val.y3.share.stayer2} of them two), and their rows carry most of the first term for the variance of
the sum.
\item \emph{Half of the workers.} A random half of the workers of the six-year panel:
\R{val.half.n.person_years} person-years of \R{val.half.n.workers} workers at \R{val.half.n.firms} firms,
\R{val.half.n.rows} rows, with a thinner network of moves.
\end{itemize}
Each sub-panel was pruned to the set in which every mover's match can be left out, the effects fitted on it
serve as its truth, and six of the ten scenarios (1, 2, 4, 6, 8 and 10) were run on it. Every cell has 200
draws. In these runs the estimator is that of Section~\ref{sec:estimator}, with no projection of the covariance
matrix (the cases in which that section gives no standard error did not arise). A draw whose variance estimate is not positive is counted and enters the coverage as a
miss; a draw that fails altogether fails the run.

The thresholds were, for every scenario and component, a ratio as in Table~\ref{tab:sim} between
\R{val.acc.ratio_min} and \R{val.acc.ratio_max} and a coverage of the 95\% interval of at least
\R{val.acc.coverage_min}; and for every run, an average coverage of at least \R{val.acc.mean_coverage_min}, no
failed draw, no variance estimate that is not positive, and floors that bind for at most \R{val.acc.max_floored}
of movers' rows in any draw. These are tolerances chosen in advance. Meeting them does not establish a
coverage of 95\%: with 200 draws the coverage of a cell is known to about $\pm0.03$.

\begin{table}[H]
\centering\footnotesize\setlength{\tabcolsep}{4.5pt}
\caption{Validation runs: root-mean-square standard error over the true one, and coverage, by run}\label{tab:val}
\begin{tabular}{@{}lccccccc@{}}
\toprule
 & & \multicolumn{4}{c}{Ratio, smallest to largest over scenarios} & & Thresholds \\
\cmidrule(lr){3-6}
Run & Scenarios & $\sigma^2_\alpha$ & $\sigma^2_\psi$ & $\sigma_{\alpha,\psi}$ & $\sigma^2_{\alpha+\psi}$ & Coverage & met \\
\midrule
\tabinput{gen/tab_val.tex}
\bottomrule
\multicolumn{8}{@{}p{0.9\textwidth}@{}}{\rule{0pt}{2.6ex}200 draws per scenario; coverage is the smallest and
the largest over scenarios and components. Appendix~\ref{app:sim} gives every cell.}
\end{tabular}
\end{table}

Two of the three runs met every threshold; the three-year window did not (Table~\ref{tab:val}). On the
six-year panel the fresh draws reproduce Table~\ref{tab:sim}: the forty ratios lie between
\R{val.six.ratio.min} and \R{val.six.ratio.max} and differ from those of the earlier runs by at most
\R{val.six.devdiff.max}, although this run has no projection of the covariance matrix. That step was therefore
immaterial; in all \R{val.all.draws} validation draws the matrix was positive definite, with a smallest
eigenvalue of at least \R{val.all.eig.min} of the largest. Coverage on the six-year panel is between
\R{val.six.cover.min} and \R{val.six.cover.max}. With half of the workers the ratios are between
\R{val.half.ratio.min} and \R{val.half.ratio.max} and coverage is between \R{val.half.cover.min} and
\R{val.half.cover.max}. In the three-year window the ratios are between \R{val.y3.ratio.min} and
\R{val.y3.ratio.max}, and the upper threshold of \R{val.acc.ratio_max} is exceeded in two cells, both for
$\sigma^2_\psi$: \R{val.y3.N6.psi.d3} in scenario~6 and \R{val.y3.N10.psi.d3} in scenario~10. Coverage there is between
\R{val.y3.cover.min} and \R{val.y3.cover.max}. In no run did a draw fail or a variance estimate come out zero or
negative, and floors bound for at most \R{val.all.floor.max} of movers' rows in any draw.

The validation adds two things. First, the results of Section~\ref{sec:sim} are not an artefact of the draws on
which the model was chosen, or of the projection. Second, the upward discrepancy is not the same in every
design. For $\sigma^2_\psi$ the largest ratio rises from \R{val.six.psi.max} on the six-year panel to
\R{val.half.psi.max} with half of the workers and \R{val.y3.psi.max} in the three-year window, the two designs
with fewer observations behind each estimated effect. On a short panel the standard error of
$\hat\sigma^2_\psi$ can thus be too large by a sixth. Across the \R{val.all.cells} cells of the three runs the ratio
is below 0.95 in \R{val.all.n.below95} and above 1.10 in \R{val.all.n.above110}. As before, none of this says how
large the model's error is on the real panel.
